% \documentclass[aspectratio=169,notes]{beamer}
\documentclass[aspectratio=169]{beamer}
\usetheme[faculty=phil]{fibeamer}
\usepackage{polyglossia}
\setmainlanguage{english} %% main locale instead of `english`, you
%% can typeset the presentation in either Czech or Slovak,
%% respectively.
\setotherlanguages{russian} %% The additional keys allow
%%
%%   \begin{otherlanguage}{czech}   ... \end{otherlanguage}
%%   \begin{otherlanguage}{slovak}  ... \end{otherlanguage}
%%
%% These macros specify information about the presentation
\title[Artificial Intelligence Software (AIS)]{Artificial Intelligence Software, HW 8} %% that will be typeset on the
\subtitle{Gradio
\\ \  \\ \  } %% title page.
\author{Oleg Bulichev}
%% These additional packages are used within the document:
\usepackage{ragged2e}  % `\justifying` text
\usepackage{booktabs}  % Tables
\usepackage{tabularx}
\usepackage{tikz}      % Diagrams
\usetikzlibrary{calc, shapes, backgrounds}
\usetikzlibrary{decorations.pathreplacing,calligraphy,calc,graphs}
\usepackage{amsmath, amssymb}
\usepackage{url}       % `\url`s
\usepackage{listings}  % Code listings
% \usepackage{subfigure}
\usepackage{floatrow}
\usepackage{subcaption}
\usepackage{mathtools}
\usepackage{todonotes}
\usepackage{fontspec}
\usepackage{multicol}
\usepackage{pdfpages}
\usepackage{wrapfig}
\usepackage{qrcode}
\usepackage{animate}
\usepackage{booktabs}
\usepackage{multirow}
% \usepackage{graphicx}
\usepackage{colortbl}

\graphicspath{{resources/}}
\frenchspacing

\setbeamertemplate{caption}[numbered]
\usetikzlibrary{graphs}

% \usepackage[backend=biber,style=ieee,autocite=footnote]{biblatex}
% \addbibresource{biblio.bib}
% \DefineBibliographyStrings{english}{%
%   bibliography = {References},}

\newcommand{\oleg}[2][] {\todo[color=red, #1] {OLEG:\\ #2}}
\newcommand{\fbckg}[1]{\usebackgroundtemplate{\includegraphics[width=\paperwidth]{#1}}}%frame background

\usepackage[framemethod=TikZ]{mdframed}
\newcommand{\dbox}[1]{
\begin{mdframed}[roundcorner=3pt, backgroundcolor=yellow, linewidth=0]
\vspace{1mm}
{#1}
\vspace{1mm}
\end{mdframed}
}

\begin{document}
\setlength{\abovedisplayskip}{0pt}
\setlength{\belowdisplayskip}{0pt}
\setlength{\abovedisplayshortskip}{0pt}
\setlength{\belowdisplayshortskip}{0pt}

\fbckg{fibeamer/figs/title_page.png}
\frame[c]{\setcounter{framenumber}{0}
    \usebeamerfont{title}%
    \usebeamercolor[fg]{title}%
    \begin{minipage}[b][6.5\baselineskip][b]{\textwidth}%
        \textcolor{black}{\raggedright\inserttitle}
    \end{minipage}
    % \vskip-1.5\baselineskip

    \usebeamerfont{subtitle}%
    \usebeamercolor[fg]{framesubtitle}%
    \begin{minipage}[b][3\baselineskip][b]{\textwidth}
        \raggedright%
        \insertsubtitle%
    \end{minipage}
    \vskip.25\baselineskip
}
%   \frame[c]{\maketitle}

\fbckg{fibeamer/figs/common.png}

\note{\scriptsize \begin{itemize}
        \item \
    \end{itemize}}

\begin{frame}[t]{Goal of the Assignment}
\vspace*{-0.4cm}
\textbf{Short description:} Build interactive Gradio web applications covering REST API integration, UI components, and ML model deployment.

\begin{enumerate}
    \item Get familiar with the Gradio framework for building interactive web demos.
    \item Integrate external REST APIs with various input/output components.
    \item Use advanced Gradio layout components: tabs, gallery, dataframe, JSON.
    \item Build a Gradio interface for the YOLO instance segmentation model from HW 7.
\end{enumerate}

\end{frame}

\begin{frame}[t]{Task 1}
\vspace{-0.3cm}
\textbf{Short description}: Build a joke application and extend it with automatic Russian translation.

\textbf{Artifacts}: Code (\texttt{.py} files) and PDF report with screenshots for each subtask.

\textbf{Requirements}
\footnotesize
\begin{enumerate}
    \item Implement a Gradio app that fetches and displays a random joke on button press.\\
    Use API: \url{https://official-joke-api.appspot.com/random_joke}
    \item Extend the app: add a button to translate the joke into Russian.\\
    Suggested model: \texttt{facebook/nllb-200-distilled-600M}
\end{enumerate}

\end{frame}

\begin{frame}[t]{Task 2}
\vspace{-0.3cm}
\textbf{Short description}: Build a weather application and an image gallery application.

\textbf{Artifacts}: Code (\texttt{.py} files) and PDF report with screenshots for each subtask.

\textbf{Requirements}
\footnotesize
\begin{enumerate}
    \item Implement a Gradio app that returns temperature and wind speed by user-provided coordinates.
    The user selects the temperature unit: Celsius or Fahrenheit.
    Use API: \url{https://open-meteo.com/}
\end{enumerate}

\end{frame}

\begin{frame}[t]{Task 3}
\vspace{-0.3cm}
\textbf{Short description}: Implement a full CRUD GUI for a REST API using Gradio tabs.

\textbf{Artifacts}: Code (\texttt{.py} files) and PDF report with screenshots for each operation.

\textbf{Requirements}
\footnotesize
Implement a Gradio app providing GUI for \texttt{/posts} at \url{https://jsonplaceholder.typicode.com/}.
Each CRUD operation is placed on a separate tab; HTTP response body is shown in \texttt{gr.JSON}.
\begin{itemize}
    \item \textbf{Read N}: fetch first N posts $\rightarrow$ \texttt{gr.Dataframe} with fields \textit{userId, id, title, body}
    \item \textbf{Read one}: fetch post by ID $\rightarrow$ fields: \textit{userId, title, body}
    \item \textbf{Create}: create new post $\rightarrow$ fields: \textit{userId, title, body}
    \item \textbf{Update}: update post by ID $\rightarrow$ fields: \textit{userId, title, body}
    \item \textbf{Delete}: delete post by ID
\end{itemize}
Reference: \url{https://huggingface.co/learn/llm-course/en/chapter9/7}

\end{frame}

\begin{frame}[t]{Task 4}
\vspace{-0.3cm}
\textbf{Short description}: Build a Gradio interface for the YOLO instance segmentation model deployed via Triton Inference Server (with FastAPI) from HW 7.

\textbf{Artifacts}: Code (\texttt{.py} files) and PDF report with screenshots.

\textbf{Requirements}
\footnotesize
\begin{itemize}
    \item Reuse the full deployment stack from HW 7 (Triton + FastAPI service).
    \item The user uploads an image in the Gradio UI; the app sends it to the \textbf{FastAPI} endpoint from HW 7 and displays the returned annotated image with segmentation masks and class labels overlaid.
    \item Inference must go through the HW 7 service stack, not directly through the YOLO model.
\end{itemize}

\end{frame}


\fbckg{fibeamer/figs/last_page.png}
\frame[plain]{}

\end{document}